\section{Physical maps and closures of \texorpdfstring{$G$}{G}-injective tensors}%
\label{sec:peps_closure_physical_maps}

Linear maps applied on every site commute with the contraction of the torus.
This section uses them, with the averaging projector, to show that closures in
distinct pair-conjugacy classes give linearly independent vectors for regular
and semi-regular $G$-injective tensors.

\begin{definition}[Products of physical linear maps]%
    \label{def:peps_physical_product_map}
    \lean{TNLean.PEPS.physicalMapSite,
        TNLean.PEPS.physicalProductMatrix,
        TNLean.PEPS.physicalProductMap,
        TNLean.PEPS.physicalProductMap_apply,
        TNLean.PEPS.torusPhysicalMap,
        TNLean.PEPS.torusPhysicalMap_apply}
    \leanok
    Let $F$ be a rectangular physical matrix. Its product over a finite
    site set $\Lambda$ has coefficients
    \begin{align}
        (F^{\otimes\Lambda})_{\tau,\sigma}
        =\prod_{v\in\Lambda}F_{\tau(v),\sigma(v)}.
    \end{align}
    The transformed site tensor is
    $(FA)^q_{t,r,b,l}=\sum_s F_{q,s}A^s_{t,r,b,l}$.
    Input and output physical dimensions may differ.
\end{definition}

\begin{lemma}[Physical maps commute with torus contraction]%
    \label{thm:peps_torus_physical_map_contraction}
    \lean{TNLean.PEPS.siteMap_physicalMapSite,
        TNLean.PEPS.torusPhysicalMap_torusBondNetwork,
        TNLean.PEPS.torusPhysicalMap_torusGClosure}
    \leanok
    \uses{def:peps_physical_product_map, def:peps_torus_g_closure}
    Applying $F$ at every physical site composes the local site map
    with $F$ and preserves every virtual bond operator. In particular,
    \begin{align}
        F^{\otimes\Lambda}\ket{\mathcal M(A\mid g,h)}
        =\ket{\mathcal M(FA\mid g,h)}.
    \end{align}
    As the physical inverse
    operation used in \cite[Theorem~5.9]{Schuch2010PEPS}, the identity holds for arbitrary torus bond operators, with no
    invariance or injectivity hypothesis.
\end{lemma}

\begin{proof}\leanok
    Expanding the physical product map and the virtual contraction,
    interchanging the finite sums, and factoring the physical sum over sites
    replaces each site coefficient by $\sum_s F_{q,s}A^s$ while leaving
    the virtual bond coefficients unchanged.
\end{proof}

\begin{lemma}[Physical product maps and Schmidt rank]%
    \label{thm:peps_physical_product_cut_rank}
    \lean{TNLean.PEPS.physicalProductMatrix_split,
        TNLean.PEPS.physicalCutMatrix,
        TNLean.PEPS.physicalCutMatrix_physicalProductMap,
        TNLean.PEPS.rank_physicalCutMatrix_physicalProductMap_le,
        TNLean.PEPS.rank_physicalCutMatrix_eq_of_physicalProductMaps}
    \leanok
    \uses{def:peps_physical_product_map}
    For a partition $\Lambda=R\sqcup R^c$, let $C_R(\psi)$ be the
    coefficient matrix of a physical vector in the two sets of
    configuration coordinates. Then
    \begin{align}
        C_R(F^{\otimes\Lambda}\psi)
        =F^{\otimes R}C_R(\psi)(F^{\otimes R^c})^{\mathsf T}.
    \end{align}
    Consequently its rank cannot increase. In the rank argument
    used in \cite[Corollary~6.10]{Schuch2010PEPS}, if physical product maps
    send $\psi$ to $\phi$ and $\phi$ back to $\psi$, their Schmidt
    ranks agree across every partition, even when neither product map is
    invertible on the full physical space.
\end{lemma}

\begin{proof}\leanok
    Splitting a physical configuration into its restrictions to $R$ and
    $R^c$ factors the product coefficient over these restrictions,
    so the resulting two physical sums give the two matrix products.
    Since matrix multiplication cannot increase rank, applying this inequality
    to the forward and reverse actions proves equality.
\end{proof}

\begin{lemma}[Recovering site coefficients from physical composition]%
    \label{thm:peps_site_coefficients_composition}
    \lean{TNLean.PEPS.siteMap_injective,
        TNLean.PEPS.physicalMapSite_eq_iff_siteMap_comp}
    \leanok
    \uses{def:peps_g_injective, def:peps_physical_product_map}
    The four-leg coefficients are uniquely determined by the site map.
    For a physical matrix $F$ and tensors $A,B$, the coefficient identity
    $FA=B$ is equivalent to $F\mathcal P(A)=\mathcal P(B)$.
\end{lemma}

\begin{proof}\leanok
    The entry of $\mathcal P(A)$ in physical row $s$ and virtual column
    $(t,r,b,l)$ is $A^s_{t,r,b,l}$. Composition by $F$ gives exactly
    $\sum_sF_{q,s}A^s_{t,r,b,l}$.
\end{proof}

\begin{lemma}[Constancy of labels away from the torus seams]%
    \label{thm:peps_torus_nonseam_connectivity}
    \lean{TNLean.PEPS.IsTorusNonseamCompatible.eq_origin,
        TNLean.PEPS.IsTorusNonseamCompatible}
    \leanok
    \uses{def:peps_torus_regular_coordinates}
    Let $q:\mathbb Z_w\times\mathbb Z_h\to X$ be labels in any set $X$, where $w,h>0$.
    Suppose that $q_{x+1,y}=q_{x,y}$ whenever $x+1\ne0$, and
    $q_{x,y}=q_{x,y+1}$ whenever $y+1\ne0$. By the synchronization away from the closure seams used in
    \cite[Theorem~5.9]{Schuch2010PEPS}, $q_{x,y}=q_{0,0}$ at every vertex.
\end{lemma}
\begin{proof}\leanok
    Inducting along the representatives $0,\ldots,w-1$ in the horizontal direction
    and then along $0,\ldots,h-1$ in the vertical direction avoids crossing
    any closure seam; the same argument includes a circumference of one.
\end{proof}

\begin{definition}[The averaging-projector site]%
    \label{def:peps_regular_averaging_site}
    \lean{TNLean.PEPS.averagingSite,
        TNLean.PEPS.siteMap_averagingSite,
        TNLean.PEPS.averagingSite_apply,
        TNLean.PEPS.averagingSite_invariant,
        TNLean.PEPS.isGInjective_averagingSite}
    \leanok
    \uses{def:peps_g_injective}
    Let $U$ be any representation of a finite group $G$ on a virtual
    space with finite basis $V$, and write $R_q$ for its four-leg
    representation. Take the physical basis to be $V^4$ and define $A_0$ by
    \begin{align}
        (A_0)^s_{t,r,b,l}
                        & =\frac1{|G|}\sum_{q\in G}(R_q)_{s,(t,r,b,l)}, \\
        \mathcal P(A_0) & =\Pi=\frac1{|G|}\sum_{q\in G}R_q.
        \label{eq:peps_averaging_site}
    \end{align}
    Then $\Pi R_q=\Pi$ and $\Pi$ is the identity on invariant vectors,
    so $A_0$ is $G$-injective. Unitarity is not required for these
    statements \cite[Definition~5.1]{Schuch2010PEPS}.
\end{definition}

\begin{lemma}[Isometry of the regular projector site]%
    \label{thm:peps_regular_averaging_site_isometry}
    \lean{TNLean.PEPS.isGIsometric_regularAveragingSite}
    \leanok
    \uses{def:peps_regular_averaging_site, def:peps_torus_regular_coordinates}
    For four regular virtual legs, $A_0$ is $G$-isometric with
    normalization factor one.
\end{lemma}

\begin{proof}\leanok
    The unitary group average is self-adjoint and idempotent, so
    $\Pi^\dagger\Pi=\Pi$. Its invariance was established in
    Definition~\ref{def:peps_regular_averaging_site}.
\end{proof}

\begin{lemma}[Group-average expansion of the torus projector contraction]%
    \label{thm:peps_torus_projector_expansion}
    \lean{TNLean.PEPS.torusBondNetwork_single,
        TNLean.PEPS.torusBondSiteLabelsEquiv,
        TNLean.PEPS.torusBondNetwork_sum,
        TNLean.PEPS.torusBondNetwork_mul,
        TNLean.PEPS.torusDress_single,
        TNLean.PEPS.torusBondNetwork_averagingSite}
    \leanok
    \uses{def:peps_torus_regular_coordinates, def:peps_regular_averaging_site}
    Let $U$ be a finite-dimensional representation of a finite group $G$, and let
    $\Pi$ be the four-leg averaging projector. Write $e(v)$ and $n(v)$ for the
    neighbours east and north of $v$, and let $s_v=(t_v,r_v,b_v,l_v)$ denote the
    physical basis label of the projector site. For arbitrary bond operators
    $O^H_v,O^V_v$, set
    \begin{align}
        H_v(q)=U_{q_{e(v)}}O^H_vU_{q_v^{-1}},\qquad
        V_v(q)=U_{q_v}O^V_vU_{q_{n(v)}^{-1}}.
        \label{eq:peps_torus_projector_bond_operators}
    \end{align}
    The torus contraction of the projector sites has coefficient
    \begin{align}
        \Psi_\Pi(s)=|G|^{-N}\sum_{q:\mathbb Z_w\times\mathbb Z_h\to G}
        \prod_v H_v(q)_{l_{e(v)},r_v}\,V_v(q)_{t_v,b_{n(v)}},
        \qquad N=wh.
    \end{align}
    As in the expansion of the projector contraction in
    \cite[Theorem~5.9]{Schuch2010PEPS}, no semi-regularity or unitarity assumption is required for this identity.
\end{lemma}
\begin{proof}\leanok
    \uses{lem:peps_torus_bond_network_gauge}
    The bond-end indices are in bijection with the four indices read at the sites.
    For identity sites, the Kronecker deltas leave the product of the bond entries.
    Expanding the group average independently at every vertex and applying finite distributivity
    gives a sum over vertex labels and the factor $|G|^{-N}$. Absorbing each
    four-leg operator into the neighbouring bonds gives $H_v(q)$ and $V_v(q)$ as in
    (\ref{eq:peps_torus_projector_bond_operators}).
\end{proof}

\begin{lemma}[Extracting the torus closure operators]%
    \label{thm:peps_torus_projector_extraction}
    \lean{TNLean.PEPS.matrixLinearFunctional_eq_sum,
        TNLean.PEPS.torusOutputBondPairing,
        TNLean.PEPS.torusOutputBondPairing_bondProduct,
        TNLean.PEPS.torusDeltaPairing,
        TNLean.PEPS.torusSeamPairing,
        TNLean.PEPS.torusDeltaPairing_apply_rep,
        TNLean.PEPS.torusSeamPairing_apply_rep,
        TNLean.PEPS.torusMatrixEntry,
        TNLean.PEPS.torusHorizontalExtractionPairing,
        TNLean.PEPS.torusVerticalExtractionPairing,
        TNLean.PEPS.torusProjectorExtraction,
        TNLean.PEPS.torusProjectorExtraction_bondProduct,
        TNLean.PEPS.torusDeltaPairing_apply_mul_inv,
        TNLean.PEPS.torusHorizontalExtractionPairing_constant,
        TNLean.PEPS.torusVerticalExtractionPairing_constant,
        TNLean.PEPS.torusProjectorExtraction_constant_bondProduct,
        TNLean.PEPS.torusHorizontalExtractionPairing_nonseam,
        TNLean.PEPS.torusVerticalExtractionPairing_nonseam,
        TNLean.PEPS.torusProjectorExtraction_label_bondProduct,
        TNLean.PEPS.torusProjectorExtraction_sum_label_bondProduct,
        TNLean.PEPS.torusProjectorExtraction_torusGClosure,
        TNLean.PEPS.torusProjectorExtraction_normalization_pos}
    \leanok
    \uses{def:peps_semiregular, def:peps_regular_averaging_site,
        def:peps_torus_g_closure, def:peps_delta_operator}
    Let $U$ be a semi-regular representation of a finite group $G$ on
    a virtual space with finite basis $V$. Let $\Lambda$ be a rectangular
    torus with positive width and height, and let $A_0$ be the averaging-
    projector site. There is one explicit linear map
    $E:\C^{\Lambda\to V^4}\to\operatorname{Mat}_{V\times V}(\C)$,
    independent of $g,h$, such that
    \begin{align}
        E\ket{\mathcal M(A_0\mid g,h)}
        =|G|^{-|\Lambda|}\sum_{x\in G}
        U_{xhx^{-1}}\otimes U_{xgx^{-1}}.
        \label{eq:peps_projector_extraction}
    \end{align}
    Here the tensor product is written in the product basis, with the
    horizontal factor first. The coefficient $|G|^{-|\Lambda|}$ is
    strictly positive. This is an explicit contraction realizing the
    independence argument of \cite[Theorem~5.9]{Schuch2010PEPS}.
\end{lemma}

\begin{proof}\leanok
    \uses{thm:peps_semiregular_trace, thm:peps_torus_projector_expansion,
        thm:peps_torus_nonseam_connectivity}
    Put $\delta_p(M)=\tr(U_{p^{-1}}M\Delta)$ and
    $\phi=\sum_p\delta_p$. Trace duality gives
    $\delta_p(U_z)=\delta_{p,z}$ and $\phi(U_z)=1$.
    Retain the horizontal seam bond $(-1,0)$ and the vertical seam bond
    $(0,-1)$. Apply $\delta_1$ to every nonseam output bond and $\phi$
    to every other seam output bond. On each retained bond use a matrix
    entry functional. Explicitly, for fixed rows $(i,k)$ and columns
    $(j,l)$, denote these horizontal and vertical functionals by
    $\kappa^h_v$ and $\kappa^v_v$, respectively. If $E_{a,b}$ is a
    matrix unit and $\beta$ assigns four end labels to each pair of bonds,
    define
    \begin{align}
        (Ef)_{(i,k),(j,l)}
        =\sum_\beta f(\mathcal A(\beta))\prod_{v\in\Lambda}
        \kappa^h_v(E_{\beta_v^{hH},\beta_v^{hT}})
        \kappa^v_v(E_{\beta_v^{vH},\beta_v^{vT}}),
    \end{align}
    where $\mathcal A(\beta)$ is the four-leg configuration read from
    the bond ends, in top, right, bottom, left order. Finite distributivity
    evaluates this contraction separately on each bond.
    For a vertex labelling $q$, a nonseam horizontal bond gives
    $\delta_1(U_{q_{v+e_1}q_v^{-1}})=\delta_{q_{v+e_1},q_v}$;
    a nonseam vertical bond gives
    $\delta_1(U_{q_vq_{v+e_2}^{-1}})=\delta_{q_v,q_{v+e_2}}$.
    Lemma~\ref{thm:peps_torus_nonseam_connectivity} forces $q_v=x$
    at every vertex. All discarded seam factors then equal one, and
    the retained factors are $U_{xhx^{-1}}$ and $U_{xgx^{-1}}$.
    The normalized projector expansion of
    Lemma~\ref{thm:peps_torus_projector_expansion} gives the coefficient
    $|G|^{-|\Lambda|}$, giving~(\ref{eq:peps_projector_extraction}).
\end{proof}

\begin{lemma}[Closure vectors and the local inverse]%
    \label{thm:peps_regular_g_injective_closure_inverse}
    \lean{TNLean.PEPS.IsGInjective.exists_torusGClosure_averagingSite,
        TNLean.PEPS.torusPhysicalMap_averagingSite}
    \leanok
    \uses{def:peps_regular_averaging_site}
    Let $A$ be $G$-injective for any finite virtual representation $U$
    and write $T=\mathcal P(A)$. There is a local rectangular matrix
    $L$ with $LT=\Pi$ such that, for every closure pair,
    \begin{align}
        L^{\otimes\Lambda}\ket{\mathcal M(A\mid g,h)}
         & =\ket{\mathcal M(A_0\mid g,h)},
        \label{eq:peps_projector_closure_forward} \\
        T^{\otimes\Lambda}\ket{\mathcal M(A_0\mid g,h)}
         & =\ket{\mathcal M(A\mid g,h)}.
        \label{eq:peps_projector_closure_reverse}
    \end{align}
    The reverse identity requires only virtual invariance of $A$.
    These are contraction identities for the inverse argument
    of \cite[Theorem~5.9]{Schuch2010PEPS}; no unitarity or local
    isometry hypothesis is needed.
\end{lemma}

\begin{proof}\leanok
    \uses{thm:peps_site_coefficients_composition, thm:peps_g_injective_left_inverse, thm:peps_torus_physical_map_contraction}
    The local inverse and virtual invariance give $LT=\Pi$ and
    $T\Pi=T$. By Lemma~\ref{thm:peps_site_coefficients_composition},
    these are coefficient identities, and applying
    Lemma~\ref{thm:peps_torus_physical_map_contraction} at every
    physical site in both directions completes the proof. No inverse on the entire
    physical space is required.
\end{proof}

\begin{lemma}[Independent closures of semi-regular $G$-injective tensors]%
    \label{thm:peps_semiregular_g_injective_closure_independence}
    \lean{TNLean.PEPS.IsGInjective.linearIndependent_torusGClosureClass_of_isSemiRegular,
        TNLean.PEPS.IsGInjective.torusGClosure_ne_zero_of_isSemiRegular,
        TNLean.PEPS.IsGInjective.linearIndependent_torusGClosureClass_commuting_of_isSemiRegular}
    \leanok
    \uses{def:peps_semiregular}
    Let $U$ be a semi-regular representation of a finite group $G$,
    and let $A$ be $G$-injective for its four-leg virtual representation.
    On every nonempty rectangular torus, the closure vectors
    $\ket{\mathcal M(A\mid C)}$, indexed by simultaneous-conjugacy
    classes $C$ of pairs, are linearly independent. Every closure
    vector is therefore nonzero. In particular, the family indexed
    by commuting pair classes is independent.

    The commuting-family assertion proves the independence component
    of \cite[Theorem~5.9]{Schuch2010PEPS}, without a local isometry
    assumption. Independence for noncommuting classes and size-one
    tori are auxiliary extensions. No identification with the
    parent-Hamiltonian ground space is asserted
    \cite{gap:rmp_peps_quantum_double_g_isometry}.
\end{lemma}

\begin{proof}\leanok
    \uses{thm:peps_regular_g_injective_closure_inverse, thm:peps_torus_projector_extraction, thm:peps_pair_conjugacy_operator_independence}
    Choosing the same local inverse $L$ for every closure vector and denoting by $N$ the number of vertices, the explicit map $E$ gives
    \begin{align}
        E L^{\otimes\Lambda}\ket{\mathcal M(A\mid C)}
         & =|G|^{-N}\Omega_C.
        \label{eq:peps_semiregular_closure_extraction}
    \end{align}
    Thus $|G|^N E L^{\otimes\Lambda}$ sends any proposed linear
    relation among the physical vectors to a relation among the
    independent operators $\Omega_C$, whose coefficients vanish by
    Lemma~\ref{thm:peps_pair_conjugacy_operator_independence}.
    Each vector is therefore nonzero, and restriction to the commuting classes
    preserves independence.
\end{proof}

\begin{lemma}[Independent regular isometric closure vectors]%
    \label{thm:peps_regular_torus_closure_independence}
    \lean{TNLean.PEPS.IsGIsometric.torusGClosure_dotProduct_eq_zero_of_pairConjugacyClass_ne,
        TNLean.PEPS.IsGIsometric.exists_torusGClosure_dotProduct_self,
        TNLean.PEPS.IsGIsometric.torusGClosure_dotProduct_self_ne_zero,
        TNLean.PEPS.IsGIsometric.torusGClosure_ne_zero,
        TNLean.PEPS.IsGIsometric.torusGClosureClass_orthogonal,
        TNLean.PEPS.IsGIsometric.linearIndependent_torusGClosureClass,
        TNLean.PEPS.IsGIsometric.linearIndependent_torusGClosureClass_commuting}
    \leanok
    \uses{thm:peps_torus_pair_conjugacy}
    For a uniform regular $G$-isometric tensor, the vector associated
    to each pair-conjugacy class is nonzero, and distinct class vectors
    are orthogonal. Their family is therefore linearly independent.
    More precisely,
    \begin{align}
        \|\psi_{g,h}\|^2=c^N|G|^N|Z_G(g)\cap Z_G(h)|>0.
    \end{align}
    Restricting to commuting classes proves the independence assertion
    underlying \cite[Theorem~5.9]{Schuch2010PEPS} in the regular
    $G$-isometric case. The all-class assertion is an auxiliary extension:
    no ground-state claim is made for noncommuting closures. Semi-regular
    $G$-injective independence is proved in
    Lemma~\ref{thm:peps_semiregular_g_injective_closure_independence};
    ground-space identification remains separate
    \cite{gap:rmp_peps_quantum_double_g_isometry}.
\end{lemma}

\begin{proof}\leanok
    \uses{thm:peps_regular_torus_closure_overlap}
    A simultaneous intertwiner exists exactly when the two closure pairs
    are simultaneously conjugate. Thus the overlap between distinct
    classes vanishes. For a self-overlap the intertwiners form
    $Z_G(g)\cap Z_G(h)$, which contains the identity. Taking the inner
    product of a finite linear relation with each nonzero class vector
    proves independence. Independence is inherited by the commuting
    subfamily.
\end{proof}

\begin{lemma}[Independent closures of regular $G$-injective tensors]%
    \label{thm:peps_regular_g_injective_closure_independence}
    \lean{TNLean.PEPS.IsGInjective.torusGClosure_ne_zero,
        TNLean.PEPS.IsGInjective.linearIndependent_torusGClosureClass,
        TNLean.PEPS.IsGInjective.linearIndependent_torusGClosureClass_commuting}
    \leanok
    \uses{def:peps_torus_regular_coordinates,
        thm:peps_semiregular_g_injective_closure_independence}
    For a regular $G$-injective tensor on any nonempty rectangular torus,
    every closure vector is nonzero and the vectors indexed by
    distinct pair-conjugacy classes are linearly independent. In
    particular the family indexed by commuting classes is independent.
    The latter is the regular-representation independence component of
    \cite[Theorem~5.9]{Schuch2010PEPS}, with no isometry hypothesis.
    Independence for all pair classes is an auxiliary extension and
    makes no ground-space assertion. The semi-regular representation result appears in
    Lemma~\ref{thm:peps_semiregular_g_injective_closure_independence};
    ground-space identification remains separate
    \cite{gap:rmp_peps_quantum_double_g_isometry}.
\end{lemma}

\begin{proof}\leanok
    \uses{def:peps_torus_regular_coordinates,
        thm:peps_semiregular_g_injective_closure_independence}
    The native regular representation is semi-regular by its group-basis
    equivalence with the group-algebra regular representation, so
    Lemma~\ref{thm:peps_semiregular_g_injective_closure_independence} applies directly.
\end{proof}

\begin{lemma}[Schmidt ranks of regular $G$-injective closure superpositions]%
    \label{thm:peps_regular_g_injective_closure_rank_transport}
    \lean{TNLean.PEPS.IsGInjective.rank_physicalCutMatrix_torusGClosure_sum}
    \leanok
    Let $A$ be regular $G$-injective. For every finite family of closure
    pairs and every physical partition $R$, the two vectors
    \begin{align}
        \sum_s\lambda_s\ket{\mathcal M(A\mid g_s,h_s)},\qquad
        \sum_s\lambda_s\ket{\mathcal M(A_0\mid g_s,h_s)}
        \label{eq:peps_closure_superposition_pair}
    \end{align}
    have the same Schmidt rank across $R$. This is the local
    deformation step of \cite[Corollary~6.10]{Schuch2010PEPS}; no
    particular boundary-rank formula is asserted here.
\end{lemma}

\begin{proof}\leanok
    \uses{thm:peps_regular_g_injective_closure_inverse,
        thm:peps_physical_product_cut_rank}
    The same forward and reverse physical product maps act on every
    summand. By linearity they recover the superpositions
    (\ref{eq:peps_closure_superposition_pair}) from each other, and Lemma~\ref{thm:peps_physical_product_cut_rank}
    gives equality of their coefficient-matrix ranks.
\end{proof}
